% ============================================================================
% LERNBLATT: Kräfte zusammensetzen (UE 2)
% Physik Klasse 7 – Gesamtschule MV
% Farbige Kräfte, kompakt auf 1 Seite
% ============================================================================
\documentclass[a4paper,9pt]{article}

\usepackage[left=1.5cm,right=1.5cm,top=1.2cm,bottom=1.2cm]{geometry}
\usepackage[utf8]{inputenc}
\usepackage[T1]{fontenc}
\usepackage[ngerman]{babel}
\usepackage{tikz}
\usepackage{amssymb,amsmath}
\usepackage{xcolor}
\usepackage{multicol}
\usepackage{enumitem}

\renewcommand{\familydefault}{\sfdefault}

\definecolor{physik}{HTML}{2c5aa0}
\definecolor{gridcolor}{RGB}{200,200,200}
\definecolor{f1color}{HTML}{1565c0}   % Blau für F1
\definecolor{f2color}{HTML}{2e7d32}   % Grün für F2
\definecolor{frescolor}{HTML}{c62828} % Rot für Fres
\definecolor{hinweis}{RGB}{120,120,120}

\setlength{\parindent}{0pt}
\setlist[enumerate]{topsep=1pt,itemsep=0pt,parsep=0pt}
\setlist[itemize]{topsep=1pt,itemsep=0pt,parsep=0pt}
\pagestyle{empty}

% Farbige Formelzeichen
\newcommand{\Fone}{\textcolor{f1color}{F_1}}
\newcommand{\Ftwo}{\textcolor{f2color}{F_2}}
\newcommand{\Fres}{\textcolor{frescolor}{F_{\text{res}}}}

% AB-Verweis
\newcommand{\abmark}[1]{\hfill{\footnotesize\textcolor{hinweis}{$\rightarrow$ AB #1}}}

% X-Kreuz (45° gedreht) am Angriffspunkt
\newcommand{\xkreuz}[2]{%
    \draw[black, thick] (#1-0.25,#2-0.25) -- (#1+0.25,#2+0.25);
    \draw[black, thick] (#1-0.25,#2+0.25) -- (#1+0.25,#2-0.25);
}

\begin{document}

% ==================== HEADER ====================
\begin{center}
{\Large\bfseries Lernblatt: Kräfte zusammensetzen}\\[1mm]
{\small Physik Klasse 7}
\end{center}

\vspace{2mm}
\hrule
\vspace{5mm}

\begin{multicols}{2}

% ==================== KRAFTPFEIL ====================
\textbf{Der Kraftpfeil} hat drei Eigenschaften: \abmark{1}
\begin{enumerate}
    \item \textbf{Angriffspunkt} – Wo greift die Kraft an?
    \item \textbf{Richtung} – Wohin zeigt die Kraft?
    \item \textbf{Betrag} – Wie groß ist die Kraft? (Länge)
\end{enumerate}

\vspace{1mm}
\begin{center}
\begin{tikzpicture}[scale=0.5]
    \draw[gridcolor, very thin] (0,0) grid (8,3);
    \xkreuz{1}{1.5}
    \draw[f1color, thick, ->] (1,1.5) -- (7,1.5);
    \node[below, font=\scriptsize] at (1,0.8) {Angriffspunkt};
    \node[above, font=\scriptsize] at (4,1.8) {Betrag};
    \node[right, font=\scriptsize] at (7.2,1.5) {Richtung};
\end{tikzpicture}
\end{center}

\vspace{5mm}
\hrule
\vspace{6mm}

% ==================== GLEICHE RICHTUNG ====================
\textbf{Kräfte in gleicher Richtung} \abmark{2}

\vspace{2mm}
\begin{center}
\fbox{\large $\Fres = \Fone + \Ftwo$}
\end{center}

\vspace{2mm}
\textbf{Beispiel:} Zwei schieben einen Wagen.

\begin{center}
\begin{tikzpicture}[scale=0.4]
    \draw[gridcolor, very thin] (0,0) grid (12,4);
    \xkreuz{1}{2}
    \draw[f1color, thick, ->] (1,2.8) -- (5,2.8);
    \node[above, font=\scriptsize, f1color] at (3,3.1) {$F_1$ = 3 N};
    \draw[f2color, thick, ->] (1,1.2) -- (7,1.2);
    \node[below, font=\scriptsize, f2color] at (4,0.9) {$F_2$ = 4 N};
    \draw[frescolor, ultra thick, ->] (1,2) -- (11,2);
    \node[above, font=\scriptsize, frescolor] at (6,2.3) {$F_{\text{res}}$ = 7 N};
\end{tikzpicture}
\end{center}

$\Fres = 3\,\text{N} + 4\,\text{N} = 7\,\text{N}$

\vspace{4mm}
\hrule
\vspace{6mm}

% ==================== ENTGEGENGESETZT ====================
\textbf{Entgegengesetzte Kräfte} \abmark{3}

\vspace{2mm}
\begin{center}
\fbox{\large $\Fres = |\Fone - \Ftwo|$}
\end{center}

\vspace{2mm}
Richtung: In Richtung der \textbf{größeren} Kraft!

\vspace{1mm}
\textbf{Beispiel:} Tauziehen.

\begin{center}
\begin{tikzpicture}[scale=0.4]
    \draw[gridcolor, very thin] (0,0) grid (13,4);
    \xkreuz{6.5}{2}
    \draw[f1color, thick, <-] (1,2.8) -- (6.5,2.8);
    \node[above, font=\scriptsize, f1color] at (3.75,3.1) {$F_1$ = 6 N};
    \draw[f2color, thick, ->] (6.5,1.2) -- (10.5,1.2);
    \node[below, font=\scriptsize, f2color] at (8.5,0.9) {$F_2$ = 4 N};
    \draw[frescolor, ultra thick, <-] (4.5,2) -- (6.5,2);
    \node[below, font=\scriptsize, frescolor] at (5.5,1.6) {$F_{\text{res}}$ = 2 N};
\end{tikzpicture}
\end{center}

$\Fres = |6\,\text{N} - 4\,\text{N}| = 2\,\text{N}$ nach links

\columnbreak

% ==================== KRÄFTEGLEICHGEWICHT ====================
\textbf{Kräftegleichgewicht} \abmark{4}

\vspace{2mm}
Wenn $\Fone = \Ftwo$ und entgegengesetzt:

\vspace{2mm}
\begin{center}
\fbox{\large $\Fres = 0\,\text{N}$}
\end{center}

\vspace{2mm}
$\Rightarrow$ \textbf{Gleichgewicht} – nichts bewegt sich!

\textbf{Beispiele:} Armdrücken (gleich stark), Buch auf Tisch

\vspace{4mm}
\hrule
\vspace{6mm}

% ==================== TIPPS ====================
\textbf{Tipps zum Zeichnen}

\textbf{Maßstab:} 1 cm $\hat{=}$ 1 N \quad (oder: 2 Kästchen $\hat{=}$ 1 N)

\vspace{1mm}
\textbf{Schritte:}
\begin{enumerate}
    \item Angriffspunkt mit $\times$ markieren
    \item Richtung festlegen
    \item Länge im Maßstab abtragen
    \item Pfeilspitze zeichnen
    \item Kraft beschriften (z.B. $\Fone$ = 3 N)
\end{enumerate}

\vspace{2mm}
\fbox{\parbox{0.95\linewidth}{%
\textbf{Achtung:} Die Pfeilspitze zählt zur Gesamtlänge!

\vspace{1mm}
\begin{center}
\begin{tikzpicture}[scale=0.45]
    % RICHTIG
    \node[anchor=east, font=\scriptsize\bfseries, text=green!50!black] at (0,1) {$\checkmark$ RICHTIG:};
    \draw[green!50!black, thick, ->] (0.5,1) -- (4.5,1);
    \draw[gray, thin] (0.5,0.4) -- (0.5,0.6) -- (4.5,0.6) -- (4.5,0.4);
    \node[font=\tiny, gray] at (2.5,0.2) {4 cm (bis Spitze)};
    % FALSCH
    \node[anchor=east, font=\scriptsize\bfseries, text=red!70!black] at (0,-0.5) {$\times$ FALSCH:};
    \draw[red!70!black, thick] (0.5,-0.5) -- (4.5,-0.5);
    \draw[red!70!black, thick, ->] (4.5,-0.5) -- (5.2,-0.5);
    \node[font=\tiny, red!70!black] at (2.5,-1.1) {4 cm + Spitze = zu lang!};
\end{tikzpicture}
\end{center}
}}

\end{multicols}

\vspace{4mm}
\hrule
\vspace{3mm}

% ==================== ZUSAMMENFASSUNG ====================
\begin{center}
\textbf{Zusammenfassung: Formeln}
\end{center}

\vspace{2mm}

\begin{center}
\renewcommand{\arraystretch}{1.5}
\begin{tabular}{|l|c|l|}
\hline
\textbf{Situation} & \textbf{Formel} & \textbf{Beispiel} \\
\hline
Gleiche Richtung & $\Fres = \Fone + \Ftwo$ & Gemeinsam schieben \\
\hline
Entgegengesetzt & $\Fres = |\Fone - \Ftwo|$ & Tauziehen \\
\hline
Gleichgewicht & $\Fres = 0$ & Armdrücken (gleich stark) \\
\hline
\end{tabular}
\end{center}

\end{document}
